% source: erdos-872/paper_template.md
% source: erdos-872/publication_source_of_truth.md

\section{Introduction}

\subsection{The game}

Fix $n \ge 2$.  The \emph{divisibility antichain saturation game} is played on
the set $\{2,3,\dots,n\}$.  Starting from the empty set, the players
alternately choose integers subject to the rule that the chosen set remains
primitive: if $a \neq b$ have been chosen, then neither divides the other.  The
game ends when no legal move remains.  We write $\L(n)$ for the value of the
game under optimal play, where Prolonger seeks to maximize the number of moves
and Shortener seeks to minimize it.

Erd\H{o}s's original problem statement does not specify which player moves
first.  Throughout this paper we assume that \emph{Prolonger moves first}.  This
matches the public discussion on the Erd\H{o}s problems forum and every
upper-bound argument used below; see \cite{Er92c,Bloom872}.

It is convenient to split the board into its lower and upper halves:
\[
L := [2,n/2]\cap\Z,
\qquad
\U := (n/2,n]\cap\Z.
\]
Every subset of $\U$ is an antichain, so the upper half behaves like a pool of
individually legal terminal moves, while the lower half carries the divisibility
constraints that drive the game.

\subsection{Prior work}

The game appears on the Erd\H{o}s problem list as problem \#872
\cite{Er92c,Bloom872}.  It is a number-theoretic analogue of the triangle-free
saturation game of Hajnal.  For the latter game, F{\"u}redi and Seress
\cite{FuSe91} proved a lower bound of order $n\log n$, while Bir{\'o}, Horn, and
Wildstrom \cite{BPW16} proved an upper bound below $0.215n^2$.

The ambient combinatorial setting also interacts with the literature on
primitive sets and divisibility antichains.  We use the primitive-set viewpoint
of Cameron and Erd\H{o}s \cite{CaEr90}, the weighted extremal perspective from
Lichtman's proof of the Erd\H{o}s primitive set conjecture \cite{Lichtman23},
and supporting multiplicative-structure estimates from
\cite{Kucheriaviy24,MartinPomerance10}.

The public upper-bound record for the present game developed rapidly in early
2026.  Price and GPT-5.2 Pro gave the first linear upper bound,
\[
  \L(n) \le \left(\frac{23}{48}+o(1)\right)n,
\]
and introduced the terminology \emph{Prolonger} and \emph{Shortener}
\cite{PriceGPTPro26}.  Subsequent forum refinements by Adenwalla, StijnC,
natso26, Xiao\_Hu, Desmond Weisenberg, and others improved this to
\[
  \L(n) \le \left(\frac{419}{1008}+o(1)\right)n \le 0.416n
\]
\cite{Forum872}.  The results in this paper substantially strengthen that
record.

\subsection{Main results}

Our first theorem is a structural reduction from terminal game positions to a
weighted lower-half optimization problem.

% source: erdos-872/current_state.md (Shield Reduction Theorem)
\begin{theorem}[Shield Reduction]\label{thm:intro-shield}
Let $P \subseteq \U$ be the set of upper-half moves played by Prolonger, and
let $\beta(P)$ be the weighted lower-half antichain parameter defined in
\Cref{sec:shield-reduction}.  Then every terminal game position $A$ satisfies
\[
  |A| \ge |\U| - \beta(P).
\]
\end{theorem}

The reduction leads to a sharp obstruction theorem for short shield prefixes and
to our first lower bound.

% source: erdos-872/current_state.md (Polynomial shield lower bound --- Theorem A)
\begin{theorem}[Polynomial shield barrier / Theorem A]\label{thm:intro-theorem-a}
For every fixed $0 < \alpha < 1$ and every $P \subseteq \U$ with $|P|\le
n^\alpha$,
\[
  \beta(P) \ge \left(\frac12 \log\frac1\alpha + o(1)\right)n.
\]
Consequently
\[
  \L(n) \ge \left(\frac18-o(1)\right)\frac{n\log\log n}{\log n}.
\]
\end{theorem}

Our second lower bound comes from a Maker-first hypergraph capture argument,
conditional on the finite safe-edge hypothesis stated in
\Cref{prop:t2-finite-capture}.

\begin{theorem}[Conditional T2 lower bound]\label{thm:intro-t2}
Assume the safe-edge hypothesis of \Cref{prop:t2-finite-capture} for the
activation graph states and residual slot-hypergraph states used in
\Cref{app:t2-proof}.  Then for every fixed $0<\delta<1/4$ there exists
$c_\delta>0$ such that
\[
  \L(n) \ge c_\delta \frac{n(\log\log n)^2}{\log n}
\]
for all sufficiently large $n$.
\end{theorem}

On the upper-bound side we obtain three successive constants.  The first two
come from the odd-prime-prefix Shortener strategy and increasingly sharp
Bonferroni truncations; the third uses the piecewise density of the Round~15
prime sequence together with a fourth-order comparison theorem.

% source: erdos-872/current_state.md (13/36 theorem)
% source: erdos-872/researcher-08-5-16-improvement.md
% source: erdos-872/researcher-60-pro-R57-repair-PROVED-theorems-2-1-and-4-1.md
\begin{theorem}[Upper bounds]\label{thm:intro-upper-bounds}
We have
\[
  \L(n) \le \left(\frac{13}{36}+o(1)\right)n,
  \qquad
  \L(n) \le \left(\frac{5}{16}+o(1)\right)n,
\]
and, under the explicit Round~15 strategy $\sigfifteen$,
\[
  \L(n) \le \left(\frac{\Wfour}{2}+o(1)\right)n
  = (0.1897112+o(1))n < 0.19n.
\]
\end{theorem}

We also determine exactly the first-hit cover constant for the auxiliary
covering problem on the lower half.

% source: erdos-872/current_state.md (5/24 first-hit cover)
\begin{theorem}[Exact $5/24$ cover theorem]\label{thm:intro-524}
Let $\taucov(n)$ denote the minimum size of an upper-half set $H\subseteq \U$
such that every $x\in L$ divides some $u\in H$.  Then
\[
  \taucov(n) = \frac{5}{24}n + O(1).
\]
\end{theorem}

Finally, we isolate two nontrivial Prolonger classes for which Shortener does
force nearly minimal game length, and we prove several obstruction theorems
showing why broader proof programs fail.

% source: erdos-872/phase4/theorem5_disjoint_carriers_note.md
% source: erdos-872/phase4/theorem6_rank3_squarefree_note.md
% source: erdos-872/phase4/block_product_carrier_mass_note.md
% source: erdos-872/researcher-52-pro-fresh-zoom-transversal-integrality-barrier.md
% source: erdos-872/researcher-53-codex-q-shadow-covering-dichotomy.md
% source: erdos-872/phase4/r56_separator_only_barrier_note.md
\begin{theorem}[Restricted classes and barriers]\label{thm:intro-restricted}
There are explicit classes of Prolonger play---including the disjoint small
prime carrier class and the squarefree rank-$\le 3$ carrier class---for which
Shortener forces $\L(n)=O_\alpha(n/\log n)$.  On the other hand, several natural
proof strategies fail: bounded reciprocal carrier mass cannot hold uniformly,
separator-only packet closures cannot yield any $o(n)$ upper bound, and the
late packet program exhibits a genuine transversal\allowbreak-integrality obstruction,
along with a q-shadow/covering dichotomy.
\end{theorem}

Taken together, these results leave open the central asymptotic question:
\begin{center}
does $\L(n)$ remain linear in $n$, or does it satisfy $\L(n)=o(n)$?
\end{center}
The unconditional lower and upper bounds in this paper show that the answer,
whichever it is, lies between $n\log\log n/\log n$ and $0.19n$; under the
finite safe-edge hypothesis of \Cref{prop:t2-finite-capture}, the lower side
improves to $n(\log\log n)^2/\log n$.

\subsection{Organization}

In \Cref{sec:shield-reduction} we prove the shield reduction theorem.
\Cref{sec:lower-bounds} contains T1 and T2.  In
\Cref{sec:524-cover,sec:intermediate-upper} we record the exact $5/24$ cover
theorem and the intermediate upper bounds $13/36$ and $5/16$.
\Cref{sec:main-upper} proves the main upper bound below $0.19n$, incorporating
the prime-rounding bridge of \Cref{thm:prime-bridge}.
\Cref{sec:restricted-class} treats the two restricted carrier classes.
\Cref{sec:proof-class-barriers,sec:structural-negatives} collect the positive
packet dichotomy and the structural obstruction results.
\Cref{sec:formal-verification} summarizes the Lean and Aristotle artifacts, and
the appendices record the conditional T2 proof, numerical checks,
formalization metadata, and glossary information.
